\chapter{Implementace}
\label{chap:implementation}

V této kapitole popíšeme implementaci knihovny PlanG. V předchozí kapitole jsme
navrhli rozhraní, které má sjednotit práci s generátory \texttt{geng} a
\texttt{plantri}. Nyní se podíváme na to, jak je toto rozhraní napojeno na původní
C programy a jakým způsobem zpřístupňujeme grafy uživatelským callbackům.

C++ část knihovny převádí konfiguraci uloženou v objektu na argumenty původního
programu. Callbacky a konfiguraci obarvení, které nejsou součástí argumentů
příkazové řádky, předává před spuštěním generátoru do C mezivrstvy. Aktuální
graf následně zpřístupňuje pomocí view nad interními daty generátoru.

\section{Architektura knihovny}
\label{sec:implementation-architecture}

Implementace je rozdělena do několika vrstev. Nejníže zůstávají původní zdrojové
kódy generátorů. Tyto části obsahují samotné algoritmy pro generování grafů a knihovna je
nepřepisuje, s výjimkou úprav generátoru \texttt{geng}, které doplňují podporu
barevných tříd a zakotvených vrcholů.

Nad nimi je tenká C mezivrstva. Tato vrstva zpřístupňuje vybrané funkce
a globální proměnné původních generátorů a umožňuje zaregistrovat callbacky z C++
části programu.

Nad C mezivrstvou jsou implementovány C++ backendy. Každý backend uchovává konfiguraci
generování, kontroluje její základní správnost a umí ji převést na argumenty
původního generátoru. Backend také před spuštěním registruje callbacky a předává
informace o barvách vrcholů.

Nejvyšší vrstvu tvoří šablona
\texttt{Generator\textless Backend\textgreater}, která uživateli poskytuje společné
rozhraní nezávislé na konkrétním generátoru.


\section{Vrstva \texttt{Generator} a backendy}
\label{sec:implementation-generator-backends}

Šablona \texttt{Generator\textless Backend\textgreater} dědí z konkrétního backendu.
Díky tomu jsou metody backendu dostupné přímo na objektu generátoru, ale společná
vrstva může zároveň definovat operace, které mají všechny backendy stejné. Patří sem
zejména metoda \texttt{run} a metody pro registraci callbacků
\texttt{setPrune}, \texttt{setPreprune} a \texttt{setOutproc}.

Při spuštění generování metoda \texttt{run} zavolá backendovou metodu
\texttt{prepare\_args}. Metoda \texttt{prepare\_args} vytvoří seznam řetězců které
odpovídají argumentům příkazové řádky původního programu. Poté se zavolá
\texttt{apply\_runtime\_state},
která nastaví callbacky a případně nastaví konfiguraci obarvení.
Nakonec se seznam
řetězců převede na pole ukazatelů typu \texttt{char*} a předá se původní vstupní
funkci generátoru.

Konkrétní backend tedy plní dvě role. Navenek nabízí typovější C++ metody pro
nastavení parametrů (počet vrcholů, omezení stupňů nebo výstupní formát).
Uvnitř ví, jak tyto volby vyjádřit pomocí přepínačů původního programu. Díky tomu
uživatel nemusí skládat argumenty příkazové řádky ručně, a zároveň zůstává zachován
původní způsob spuštění generátorů.

\section{Propojení s \texttt{geng}}
\label{sec:implementation-geng}

Backend pro \texttt{geng} uchovává parametry odpovídající přepínačům původního
programu. Například počet vrcholů, rozsah počtu hran, požadavky na souvislost,
omezení minimálního a maximálního stupně, volba výstupního formátu nebo rozdělení
výpočtu na části. Metoda \texttt{prepare\_args} tyto hodnoty převede na argumenty,
které by uživatel jinak zadal na příkazové řádce.

Samotné spuštění probíhá přes přejmenovanou vstupní funkci původního programu.
Kód \texttt{geng} je při sestavení upraven tak, aby jeho hlavní funkce nebyla
\texttt{main}, ale funkce volatelná z knihovny. Zároveň je generátor sestaven s
makry \texttt{OUTPROC}, \texttt{PRUNE} a \texttt{PREPRUNE}, která směřují do C
mezivrstvy. Díky těmto makrům se výpis nalezených grafů a volání prořezávacích testů
přesměrují z původního programu do funkcí knihovny.

Volání z \texttt{geng} projdou přes C mezivrstvu do trampoline funkcí v C++
části. C++ část uchovává uživatelské callbacky a při jejich volání vytvoří
\texttt{geng::GraphView}. Ten obsahuje ukazatel na graf v reprezentaci nauty,
počet vrcholů aktuálního grafu a cílový počet vrcholů.


Důležitý rozdíl oproti původnímu \texttt{PRUNE} mechanismu je v tom, jakou
semantiku nabízíme uživateli. V původním \texttt{geng} se \texttt{PRUNE} volá i
pro mezilehlé grafy během konstrukce. V rozhraní PlanG chceme, aby
\texttt{setPrune} filtroval hotové grafy, zatímco \texttt{setPreprune} sloužil k
prořezávání rozpracovaných větví. C++ napojení proto u callbacku \texttt{setPrune}
nejprve ověří, zda aktuální počet vrcholů odpovídá cílovému počtu vrcholů. 
Pokud podmínka není splněna, uživatelský callback se nezavolá.

\section{Propojení s \texttt{plantri}}
\label{sec:implementation-plantri}

Backend pro \texttt{plantri} pracuje podobným způsobem, ale napojuje se na jiné
vnitřní mechanismy původního generátoru. Konfigurace obsahuje například volbu třídy
generovaných grafů, omezení minimálního stupně a souvislosti, volbu výstupního
formátu, generování duálních grafů nebo rozdělení výpočtu. Metoda
\texttt{prepare\_args} z těchto voleb sestaví argumenty pro původní program
\texttt{plantri}.

Callbacky v \texttt{plantri} jsou navázány na makra \texttt{FILTER} a
\texttt{PRE\_FILTER}. C mezivrstva je přesměruje do trampoline funkcí v C++
části, které pak volají callbacky nastavené backendem. Callback \texttt{setPrune} odpovídá
mechanismu \texttt{FILTER} a pracuje s kandidátem na výstupní graf. Přes stejnou trampoline funkci obsluhujeme také \texttt{setOutproc}, který se
zavolá při zpracování nalezeného grafu. Callback \texttt{setPreprune} odpovídá mechanismu
\texttt{PRE\_FILTER} a může být volán nad rozpracovaným grafem během rekurzivní
konstrukce.

Protože \texttt{plantri} drží aktuální graf v globálních proměnných, C mezivrstva
zpřístupňuje funkce pro čtení těchto hodnot. C++ vrstva z nich sestaví view pro
backend \texttt{plantri}. Do něj se předávají ukazatele na hrany, pole stupňů vrcholů, počet vrcholů,
počet orientovaných hran a cílový počet vrcholů.


\section{Implementace \texttt{GraphView}}
\label{sec:implementation-graph-view}

Pro každý backend je implementován samostatný typ \texttt{GraphView}, protože
oba generátory ukládají graf odlišným způsobem. Typ \texttt{geng::GraphView}
uchovává ukazatel na graf v bitové reprezentaci nauty a počet jeho vrcholů.
Při průchodu grafem iterátory procházejí sousedy a hrany přímo v této reprezentaci.

Typ \texttt{plantri::GraphView} pracuje s poli stupňů a s orientovanými
hranami původního rovinného vložení. Při procházení hran využívá odkazy mezi
orientovanými hranami. Počítá také s hodnotou \texttt{missing\_vertex},
protože některé vnitřní reprezentace \texttt{plantri} mohou přeskočit jeden
index vrcholu.

\subsection{Integrace s Boost Graph Library}

Aby bylo možné nad objektem \texttt{GraphView} spouštět algoritmy z knihovny
Boost Graph Library, implementujeme pro oba backendy rozhraní požadované
grafovými koncepty této knihovny. Oba typy \texttt{GraphView} poskytují typy používané šablonou
\texttt{boost::graph\_traits} \footnote{\url{https://www.boost.org/doc/libs/latest/libs/graph/doc/graph_traits.html}}, 
například typ vrcholu, typ hrany, iterátory
a podporované způsoby průchodu grafem.

Pro oba backendy dále implementujeme free funkce požadované grafovými koncepty
Boost.Graph ~\cite{boostGraphConcepts}, 
například \texttt{vertices}, \texttt{edges},
\texttt{num\_vertices}, \texttt{out\_edges},
\texttt{adjacent\_vertices}, \texttt{source} a \texttt{target}.
Tyto operace pracují přímo s daty zpřístupněnými příslušným
view a nevytvářejí kopii grafu.

Součástí integrace je také mapa indexů vrcholů. U backendu \texttt{geng}
používáme identické mapování, protože vrcholy jsou číslovány souvisle od nuly.
U backendu \texttt{plantri} vlastní mapa indexů přeskočí případný chybějící
vrchol (\texttt{missing\_vertex}). Algoritmy z Boost Graph Library tak mohou pracovat s oběma backendy
stejným způsobem.

\section{UML diagram třid}
\label{sec:uml-diagram}

\section{Obarvené a zakotvené vrcholy v \texttt{geng}}
\label{sec:implementation-coloured-geng}




\section{Callbacky a přechod mezi C a C++}
\label{sec:implementation-callbacks}

\section{Výstupní formáty}
\label{sec:implementation-output-formats}
